% Pursuing Stacks (cotes 134-2 to 134-8): archive cote 134-3, site page 11 (PDF page 12), typescript page 11.
% typescript. Draft compared: ps2.tex:470-531.
% AI-transcribed by Claude Opus 5.5 (claude-opus-5-5, adaptive thinking, effort high), Message Batches API, batch ps1 (msgbatch_01Rv5hCaNTtUt4JHx1yqoHPL), 2026-11-12; not yet checked by eye.
\scanpage{134-3/11}
\mspage{11}
a large part of the development of the theory of Gr-stacks.  A systematic
study of particular sets of homotopies is closely connected \inserted{of course} to an inves-%
\hand{(}tigation into irredundancy conditions which can be figured out for a
\inserted{general} coherator $\underline{C}$ .  This is indeed an interesting topic,  but I decided not
to get involved in this,  unless I am really forced to!

The basic notion which has been peeling out in the reflection (L)
is of course the notion of a coherator.  Concerning terminology,  it oc-%
cured to me that the dual category $\underline{B}_\infty$ to $\underline{C}_\infty$ is more suggestive in some
cases,  for instance because of the topological interpretation attached
to its objects,  and (more technically)  because of formal analogy of
the role of this category,  for developping homotopical algebra,  with
\margin{(11)} the category of the standar\over{\unsure{t}}{d} (ordered) simplices\hand{(\textsuperscript{11})}.Both mutually dual
objects $\underline{B}_\infty$ , $\underline{C}_\infty$ seem to me to merit a name,  I suggest to call them
respectively \uline{left} and \uline{right} \uline{coherat ors},  or simply \uline{coherators} of course
when for a while it is understood on which side of the mirror we \unsure{are}
playing the game.

O\over{n}{\unsure{n}}e last comment still before taking off for a heuristic voyage
\over{a}{o}f discovery of stacks\hand{!} I just had a glance at Benabou\hand{'}s expos\hand{\'e} in 1967
of \over{W}{w}hat he calls ``bicategories''  (Springer Lecture N\hand{\unsure{o}}otes n\textsuperscript{o} 47 ,  p. 1
- 77).  These are non else,  it \unsure{appears}\hand{,} \deleted{\unsure{than}} than \deleted{(}\uline{non associative}
2-categories,  namely 2-stacks  in the terminology I am proposing
(but not 2-Gr-stacks  -  namely \unsure{it|at} is a particular case of a general
notion of oo-stack which  has still to be developped).  The most interes-%
\hand{-}ting feature of this exposé,  it seems to me,  is the systematic refe-%
rence to topological intuition as \inserted{a} support for a comprehension of an
algebraic formalism,  notably \inserted{of} the structure of various diagrams. His
terminology,  referring to elements of $F_o,F_1,F_2$ respectively as 0-cells,
1-cells and 2-cells,  is quite suggestive of an idea of \deleted{\unsure{geometric}} \inserted{topological} reali-%
zation of a 2-stack  -  it is not clear from this exposé \unsure{whether|whather} Benabou
has worked out this idea,  nor whether he has made a connection with
Quillen's ideas on axiomatics of homotopy theory,  which appeared the
same year in the same series\hand{(\textsuperscript{14})}. In any case,  in the last section of his
expos\hand{\'e},  he deals with his bicategories formally as with topological
spaces,  much in the same spirit as \hand{(}\inserted{the one} I was contemplating since \xout{1975xxxxx}
around 1975,  and which is now motivating the present notes.  While
there is no mention of Benabou\hand{'}s ideas in my letters\unsure{\deleted{t}} to Larry Breen,
it is quite possible that on the unconscious level,  the little I had
heard of his approach on one or two casual occasions,  had entered into
reaction with my own intuitions\hand{,} coming \inserted{mainly} from geometry and \unsure{hohomological|cohomological}
algebra,  and finally resulted in the program outlined in those letters.
